% Additional simulation results: material that the thesis does not print (see Appendix A of the thesis).
% LaTeX source, written as a chapter of the thesis. Tables it inputs are in ../tables/ (printed in the thesis)
% and ../repo-tables/ (cited from it); citation keys are those of the thesis bibliography. References to
% thesis sections, figures, tables and equations are given as their numbers in the thesis.
\chapter{Additional simulation results}\label{app:sim}

This appendix holds the simulation results that Chapter~5 uses but does not print. The full table for resampled residuals, the delays and the false-alarm figure for the independent threshold are in the repository (Appendix~A); the computing times are in Table~A.2 of the thesis.

\paragraph{Other levels of autocorrelation.} Table~\ref{tab:sim-appendix} repeats the main false-alarm and detection results at weaker ($\phi=0$) and stronger ($\phi=0.6$) autocorrelation. The pattern is the same as at $\phi=0.4$, with one exception: at $\phi=0$, the median of four and Holt-Winters exceed the false-alarm target by far without holidays, for the reason given in Section~5.1.

% written by scripts/tables.py from results/; do not edit
\begin{table}[ht]\centering\footnotesize\setlength{\tabcolsep}{6.5pt}
\caption{Simulation at $\phi=0$ and $\phi=0.6$ gives the same picture as at $\phi=0.4$, except that at $\phi=0$ the median of four and Holt-Winters exceed the false-alarm target by more; with holidays at their real size the problem is smaller but the same. FA: chance (\%) of a false alarm in a year, AR(1) threshold; Det.: share (\%) of 2 SD drops detected within 60 days, matched 5\% rate. 10{,}000 runs per cell.}\label{tab:sim-appendix}
\begin{tabular}{@{}lrrrrrrrrrr@{}}
\toprule
 & \multicolumn{2}{c}{CUSUM} & \multicolumn{2}{c}{EWMA} & \multicolumn{2}{c}{FOCuS} & \multicolumn{2}{c}{FOCuS-AR} & \multicolumn{2}{c}{CPM} \\
\cmidrule(lr){2-3}\cmidrule(lr){4-5}\cmidrule(lr){6-7}\cmidrule(lr){8-9}\cmidrule(lr){10-11}
Forecast & FA & Det. & FA & Det. & FA & Det. & FA & Det. & FA & Det. \\
\midrule
\multicolumn{11}{@{}l}{\emph{$\phi=0$, no holidays; 2 SD drop $\approx$ 9\%}} \\*
\quad True pattern & 4 & 100 & 4 & 100 & 4 & 100 & 4 & 100 & 4 & 100 \\
\quad Same-weekday difference & 1 & 41 & 0 & 45 & 4 & 27 & 4 & 26 & 1 & 40 \\
\quad Median of four & 50 & 52 & 40 & 47 & 38 & 55 & 33 & 54 & 79 & 55 \\
\quad Holt-Winters & 52 & 47 & 41 & 35 & 41 & 50 & 36 & 48 & 83 & 49 \\
\quad Yearly memory & 4 & 96 & 4 & 91 & 4 & 94 & 4 & 92 & 4 & 91 \\
\quad Raw (no forecast) & 40 & 6 & 26 & 33 & 47 & 5 & 50 & 5 & 0 & 60 \\
\addlinespace
\multicolumn{11}{@{}l}{\emph{$\phi=0$, holidays in the data; 2 SD drop $\approx$ 9\%}} \\*
\quad True pattern & 4 & 100 & 4 & 100 & 4 & 100 & 4 & 100 & 4 & 100 \\
\quad Same-weekday difference & 100 & 0 & 100 & 0 & 100 & 0 & 100 & 0 & 2 & 1 \\
\quad Median of four & 100 & 0 & 100 & 0 & 100 & 0 & 100 & 0 & 19 & 8 \\
\quad Holt-Winters & 100 & 0 & 100 & 0 & 100 & 0 & 100 & 0 & 77 & 8 \\
\quad Yearly memory & 100 & 10 & 100 & 4 & 100 & 1 & 100 & 0 & 0 & 41 \\
\quad Raw (no forecast) & 0 & 0 & 99 & 0 & 9 & 0 & 53 & 0 & 25 & 0 \\
\addlinespace
\multicolumn{11}{@{}l}{\emph{$\phi=0.6$, no holidays; 2 SD drop $\approx$ 9\%}} \\*
\quad True pattern & 6 & 100 & 6 & 98 & 6 & 100 & 6 & 98 & 5 & 99 \\
\quad Same-weekday difference & 0 & 10 & 0 & 11 & 1 & 7 & 4 & 5 & 0 & 8 \\
\quad Median of four & 7 & 22 & 8 & 17 & 7 & 21 & 7 & 15 & 7 & 21 \\
\quad Holt-Winters & 6 & 24 & 6 & 15 & 5 & 22 & 6 & 15 & 7 & 25 \\
\quad Yearly memory & 1 & 49 & 2 & 29 & 3 & 34 & 4 & 21 & 1 & 35 \\
\quad Raw (no forecast) & 43 & 4 & 38 & 26 & 47 & 4 & 50 & 5 & 0 & 41 \\
\addlinespace
\multicolumn{11}{@{}l}{\emph{$\phi=0.6$, holidays in the data; 2 SD drop $\approx$ 9\%}} \\*
\quad True pattern & 6 & 100 & 6 & 97 & 6 & 100 & 5 & 99 & 6 & 99 \\
\quad Same-weekday difference & 100 & 0 & 100 & 0 & 100 & 0 & 100 & 0 & 1 & 1 \\
\quad Median of four & 100 & 0 & 100 & 0 & 100 & 0 & 100 & 0 & 14 & 2 \\
\quad Holt-Winters & 100 & 0 & 100 & 0 & 100 & 0 & 100 & 0 & 59 & 5 \\
\quad Yearly memory & 98 & 9 & 100 & 3 & 100 & 0 & 100 & 0 & 2 & 13 \\
\quad Raw (no forecast) & 0 & 0 & 98 & 0 & 10 & 0 & 55 & 0 & 30 & 0 \\
\addlinespace
\multicolumn{11}{@{}l}{\emph{$\phi=0.4$, holidays at their real size (noise SD 0.097); 2 SD drop $\approx$ 18\%}} \\*
\quad True pattern & 6 & 100 & 5 & 100 & 6 & 100 & 6 & 100 & 6 & 100 \\
\quad Same-weekday difference & 62 & 0 & 97 & 0 & 98 & 0 & 100 & 0 & 0 & 1 \\
\quad Median of four & 87 & 1 & 100 & 0 & 100 & 0 & 100 & 0 & 3 & 21 \\
\quad Holt-Winters & 93 & 0 & 100 & 0 & 100 & 0 & 100 & 0 & 5 & 29 \\
\quad Yearly memory & 53 & 23 & 92 & 8 & 99 & 3 & 100 & 1 & 1 & 36 \\
\quad Raw (no forecast) & 1 & 2 & 21 & 0 & 1 & 0 & 2 & 0 & 4 & 0 \\
\bottomrule
\end{tabular}
\end{table}


The same drift explains the one exception at other values of $\phi$. At $\phi=0$, the median of four and Holt-Winters give a false alarm in 33 to 83\% of years without holidays (Table~\ref{tab:sim-appendix}). There the slow drift is the main autocorrelation left, which a lag-one estimate does not see; at $\phi=0.6$ the rates fall to at most 9\%. With the independent threshold almost every combination fails at $\phi=0.4$ (figure in the repository, Appendix~A). Without any forecast the rates are erratic, 0 to 98\%.

\paragraph{Variance detectors and false alarms.} Without holidays at $\phi=0.4$ the variance detectors stay close to the target with Holt-Winters and the yearly memory (4.6 to 5.8\%) and slightly above it with the two simple forecasts (6.0 to 8.6\%); like the uncapped mean detectors, they alarm in every year once holidays are in the data.

\paragraph{The Holt-Winters level weight.} A slower level ($\alpha=0.01$, not 0.05) lets capped CUSUM find 82\% of four-SD drops with holidays instead of 2\%, but the AR(1) threshold then gives CUSUM false alarms in 55\% of years without holidays; $\alpha=0.10$ meets the target but finds no more, and a frozen level alarms every year.

\paragraph{Share of variation in false alarms.} In the split of Figure~5.3, the detector explains only 2\% of the variation in false alarms without holidays. The threshold rule's share grows with the autocorrelation, from 2\% without autocorrelation to 82\% at $\phi=0.6$. With resampled noise and CPM added, the detector explains 48\% of the variation in false alarms and the forecast 1\%. For detection with holidays, the capped detectors in place of CPM leave the forecast in front: with CUSUM, EWMA, FOCuS and their capped versions, the forecast explains 63 to 76\% of the variation and the detector 6 to 31\% across Gaussian noise with simulated and with real-scale holidays, $t_5$ noise and resampled residuals (\texttt{results/extra/shares\_detsets.csv}).

\paragraph{Tuning and variance changes.} Table~\ref{tab:tuning} gives the tuning of CUSUM and EWMA, which shows how much of EWMA's gap to CUSUM comes from its weight $\lambda$. Table~\ref{tab:variance} gives the detection of variance changes behind Section~5.3.

% written by scripts/tables.py from results/; do not edit
\begin{table}[ht]\centering\footnotesize\setlength{\tabcolsep}{5pt}
\caption{A smaller smoothing constant closes most of the gap between EWMA and CUSUM for small drops (1 SD, true pattern removed: EWMA 66\% with $\lambda=0.1$ and 84\% with $\lambda=0.05$, against 89\% for CUSUM). Detection within 60 days at a matched 5\% rate and false alarms with the AR(1) threshold; $\phi=0.4$, no holidays; 10{,}000 runs per cell. Our settings: $k=0.5$, $\lambda=0.1$.}\label{tab:tuning}
\begin{tabular}{@{}lrrrrrr@{}}
\toprule
 & \multicolumn{3}{c}{CUSUM} & \multicolumn{3}{c}{EWMA} \\
\cmidrule(lr){2-4}\cmidrule(lr){5-7}
 & $k=0.25$ & $k=0.5$ & $k=1$ & $\lambda=0.05$ & $\lambda=0.1$ & $\lambda=0.2$ \\
\midrule
\multicolumn{7}{@{}l}{\emph{True pattern removed}} \\*
\quad False alarms, AR(1) (\%) & 5.5 & 5.9 & 5.7 & 5.8 & 5.8 & 5.6 \\
\quad 1 SD drop ($\approx$5\%), detected & 93 & 89 & 57 & 84 & 66 & 39 \\
\quad 2 SD drop ($\approx$9\%), detected & 100 & 100 & 100 & 100 & 100 & 99 \\
\quad 4 SD drop ($\approx$18\%), detected & 100 & 100 & 100 & 100 & 100 & 100 \\
\addlinespace
\multicolumn{7}{@{}l}{\emph{Holt-Winters}} \\*
\quad False alarms, AR(1) (\%) & 14.5 & 17.8 & 9.3 & 16.0 & 12.7 & 8.5 \\
\quad 1 SD drop ($\approx$5\%), detected & 3 & 3 & 3 & 2 & 2 & 2 \\
\quad 2 SD drop ($\approx$9\%), detected & 33 & 32 & 30 & 23 & 21 & 18 \\
\quad 4 SD drop ($\approx$18\%), detected & 99 & 100 & 100 & 96 & 97 & 95 \\
\addlinespace
\multicolumn{7}{@{}l}{\emph{Yearly memory}} \\*
\quad False alarms, AR(1) (\%) & 0.4 & 2.2 & 4.5 & 1.4 & 2.8 & 4.3 \\
\quad 1 SD drop ($\approx$5\%), detected & 18 & 11 & 6 & 12 & 8 & 5 \\
\quad 2 SD drop ($\approx$9\%), detected & 86 & 72 & 49 & 70 & 54 & 35 \\
\quad 4 SD drop ($\approx$18\%), detected & 100 & 100 & 100 & 100 & 100 & 99 \\
\bottomrule
\end{tabular}
\end{table}


% written by scripts/tables.py from results/; do not edit
\begin{table}[t]\centering\footnotesize\setlength{\tabcolsep}{4.2pt}
\caption{Detectors built for spread find a doubling of the variance 1.6 to 2.8 times as often as CUSUM, EWMA and FOCuS, and 1.3 to 1.7 times FOCuS-AR; under $t_5$ noise the advantage shrinks. Detection (\%), $\phi=0.4$, matched 5\% rate; after the yearly memory with holidays every detector finds at most 1.5\%.}\label{tab:variance}
\begin{tabular}{@{}lrrrrrrr@{}}
\toprule
 & \multicolumn{5}{c}{Holidays removed exactly} & \multicolumn{2}{c}{\begin{tabular}[b]{@{}c@{}}No holidays,\\variance $\times2$\end{tabular}} \\
\cmidrule(lr){2-6}\cmidrule(lr){7-8}
Detector & $\times1.5$ & $\times2$ & $\times3$ & \begin{tabular}[b]{@{}c@{}}$\times3$, within\\14 days\end{tabular} & $\times2$, $t_5$ & \begin{tabular}[b]{@{}c@{}}True\\pattern\end{tabular} & \begin{tabular}[b]{@{}c@{}}Yearly\\memory\end{tabular} \\
\midrule
Variance CUSUM & 10 & 31 & 75 & 11 & 12 & 30 & 24 \\
Variance EWMA & 8 & 25 & 65 & 22 & 6 & 23 & 18 \\
Variance FOCuS & 10 & 31 & 75 & 16 & 12 & 30 & 24 \\
\addlinespace
CUSUM (mean) & 4 & 11 & 28 & 4 & 8 & 11 & 7 \\
EWMA (mean) & 5 & 12 & 30 & 7 & 7 & 12 & 7 \\
FOCuS (mean) & 5 & 15 & 40 & 11 & 5 & 14 & 10 \\
FOCuS-AR (mean) & 6 & 20 & 54 & 21 & 4 & 18 & 15 \\
CPM (mean) & 2 & 3 & 4 & 0 & 2 & 3 & 2 \\
\bottomrule
\end{tabular}
\end{table}


\paragraph{Offline benchmarks.} PELT, Binary Segmentation and WBS were run with penalties calibrated to a 5\% chance of a spurious change in a normal year. On the true-pattern series they located a one-SD drop within seven days in 63 to 68\% of runs. They located a two-SD drop in 95 to 96\%. A doubling of the variance was located in at most 3.5\% of runs at $\phi = 0.4$, because their cost function looks for changes in mean. On the Holt-Winters forecast errors they located only 2 to 4\% of one-SD drops, and none when holidays were in the data. Their median location error there was 12 to 21 days. The standardisation uses the monitor's running scale, which a large step inflates for a few weeks; with one fixed scale from the training year, one-SD detection is 67 to 70\% and the spurious changes after a 4~SD step fall from 1.8 to 2.2 per run to 0.04 to 0.09.
